% !PRINTED 412
\weberpage{412}
Now if $\mu < n$, it cannot be divisible at once by $a$ and by $b$,
and hence $a$ and $b$ cannot both be the least positive exponents of the
powers of $\alpha$ and $\beta$ that become equal to 1; that is to say, if $r$ is not
a primitive $n$th root of unity, then $\alpha$ and $\beta$ are not at the same time
primitive $a$th and $b$th roots of unity; or, if
$\alpha$ and $\beta$ are primitive $a$th and $b$th roots of unity,
then their product $r$ is a primitive $n$th root of unity. On the other
hand it is clear that if $\alpha$ or $\beta$ is at the same time a root of unity of lower
degree than $a$ or $b$, then $r$ too will be a root of unity of degree lower than
the $n$th.

If $n$ decomposes into several factors $a, b, c \ldots$, of which every two are
relatively prime to one another, and if $\alpha, \beta, \gamma \ldots$ are roots of unity of the degrees
$a, b, c \ldots$, and one sets

\weq{3}{r = \alpha\beta\gamma \ldots,}
then one concludes, by repeated application of the preceding theorems, that in (3) all
the $n$th roots of unity, and each only once, are contained, and further that
$r$ is a \gesperrt{primitive} $n$th root of unity then and only then, when
$\alpha, \beta, \gamma \ldots$ are primitive roots of unity of the degrees $a, b, c \ldots$.

If we now denote the number of primitive $n$th roots of unity by
$\varphi(n)$, it follows from what has here been proved that

\weq{4}{\varphi(n) = \varphi(a)\varphi(b)\varphi(c) \ldots,}
if

\[
n = abc \ldots
\]
and $a, b, c \ldots$ are numbers which, two by two, are relatively prime to one another.

Now every number $n$ can be decomposed in one and only one way into a product of
prime powers

\[
n = p^\pi p_1^{\pi_1} p_2^{\pi_2} \ldots,
\]
in which $p, p_1, p_2 \ldots$ are distinct primes and $\pi, \pi_1, \pi_2 \ldots$
are positive exponents, so that from the formula (4) it follows:

\weq{5}{\varphi(n) = \varphi(p^\pi)\, \varphi(p_1^{\pi_1})\, \varphi(p_2^{\pi_2}) \ldots,}
and that it therefore only remains to decide whether and how many primitive
roots of unity of the degree $p^\pi$ exist.

But this question is very simple to decide. For if $\varrho$ is a
root of unity of the degree $p^\pi$, then the lowest degree to which $\varrho$ belongs as a
root of unity is a divisor of $p^\pi$, hence a power of $p$, and if it is
therefore not equal to
%CONT>>
